% Abstract -- this file is \input INSIDE the \begin{abstract}...\end{abstract}
% environment in main.tex, so DO NOT add \begin{abstract} here.
% Do NOT cite references in the abstract.
%
% NUMBERS: computed from evaluation_results (2).csv (2026-07-21 snapshot).
%   Scope: 3 attack scenarios (tool / memory / prompt-injection-MMLU)
%          x 4 topologies (random, chain, tree, star) x 55-60 dialogues,
%          backbone = Llama-3.1-70B-Instruct (AWQ-INT4). n = 12 runs per model.
%   ASR @ final turn, overall : undefended 20.8% | baseline 18.0% | T-GAT 11.6%
%   ASR @ final turn, tool    : undefended 25.6% | baseline 19.1% | T-GAT  3.4%
%   Detection recall @ first post-injection turn : baseline 73.9% -> T-GAT 81.8%
%   (All figures are ABSOLUTE percentages -- keep units consistent; do not mix in
%    percentage-point deltas or relative reductions here.)
%
% TODO(backbones): only ONE backbone measured so far; update the wording once
%   GPT / Claude / Gemini runs land.
% TODO(csqa): PI-CSQA excluded from headline numbers (single topology + accuracy
%   floor ~0.48); re-include after fixing the loader/labels.
% TODO(continuous baseline): no continuous-attack control run yet, so we claim
%   "delays detection" (measured) rather than "degrades vs continuous".
% NOTE(experiments): the escalation attack also drives multi-agent task accuracy
%   from 93.9% to 74.4% over a dialogue (-19.5 pts) -- report that in the
%   Experiments section, not here (different unit).

LLM-based multi-agent systems (MAS) achieve strong performance on complex
collaborative tasks, but their inter-agent communication opens an attack surface
in which a few compromised agents propagate malicious behavior across the system.
State-of-the-art defenses cast attacker detection as node classification on a
multi-agent \emph{utterance graph}, yet they compress each agent's multi-turn
dialogue into a single mean-pooled embedding before reasoning over the graph.
We identify this temporal compression as a critical blind spot: because detection
relies on an averaged representation, an attacker that starts benign and escalates
its intent only gradually keeps its mean embedding close to that of honest agents,
and thereby delays detection until after the damage has spread. Building on this
insight, we make two contributions. First, we introduce an \emph{escalation
attack} that saturates the averaged feature with benign early turns: it lifts the
attack success rate of the system to \textbf{20.8\%}, and the state-of-the-art
mean-pooled defense brings it down to only \textbf{18.0\%}. Second, we propose
\textbf{T-GAT}, a temporal-aware detector that replaces mean-pooling with a
temporal edge encoder fusing self-attention, temporal convolution, and cross-turn
statistics (variance and step-wise change); by modeling \emph{how} agents change
rather than only \emph{what} they say on average, T-GAT defends against both
continuous and escalation attacks. Across three attack scenarios (tool calling,
memory poisoning, and prompt injection) and four communication topologies, T-GAT
raises first-turn detection recall from \textbf{73.9\%} to \textbf{81.8\%} and
lowers the attack success rate to \textbf{11.6\%}. On the strongest tool-calling
attack, where the baseline recovers only from \textbf{25.6\%} to \textbf{19.1\%},
T-GAT reaches \textbf{3.4\%}.
